% ============================================================================
% NEW MATERIAL, October 6, 2026 (agent draft; not in JMP_DS_draft).
% Motivating-evidence section for the oct05 build. It reorganizes, without
% editing them, sections/02_empirical_evidence.tex (housing around the first
% birth, Sept 24), the fertility facts of 02b, and the "Who holds the large
% homes" subsection of the Oct 4 sandbox paper
% (tmp/jmp_paper_sandbox_20261004/paper/sections/02_empirical_evidence.tex,
% reviewed there against primary sources). Calibration data and targets now
% live in Section 4. Register: Boar, Gorea and Midrigan, Section 2.
% ============================================================================
\providecommand{\mocktodo}[1]{\textcolor{red}{[TODO: #1]}}

\section{Motivating Evidence}
\label{sec:empirical-evidence}

This section documents four facts that shape the model. First, children
arrive early in adult life, and earliest in lower-income families. Second,
large homes are mostly owned. Third, most large owner-occupied homes are held
by older households without children at home, while young families hold few.
Fourth, households enlarge their dwelling and move into ownership within a few
years of the first birth, and then move less. Together these facts describe
how a young family obtains space for children: it needs a larger home at the
start of adult life, when its savings are smallest, and the large homes are
mostly sold rather than rented and mostly held by an older generation. None
of the facts isolates an exogenous change in fertility, and I do not use them
to argue that housing costs lower fertility. Appendix~\ref{app:empirical-data}
describes the data.

\subsection{Children arrive early}

First births are concentrated in the mid twenties. In the 2003--2006 natality
records the mean age of mothers at first birth is 26.0 years, and a woman aged
25 in the 2004 and 2006 Current Population Survey (CPS) has had 0.81 children.
% [DATA] NCHS 2003-06, 25.976; CPS June 2004/06 women aged 25, 0.810 (capped at 3).
Early childbearing is concentrated in lower-income families.
Table~\ref{tab:emp-fertility-income} groups women in the 2024 CPS fertility
supplement into thirds of family income. Among women aged 24 to 26, those in
the bottom third have had 0.60 children and those in the top third 0.23. By
ages 40 to 44 the gradient has nearly disappeared: 1.78 children in the bottom
third and 1.73 in the top. Lower-income women do not have many more children
over their lives; they have them earlier, at the ages at which households
have the least wealth. The fact describes who has children early, not how
income affects fertility: young women who live with their parents report
their parents' income, and a birth itself lowers a mother's earnings.
Restricting the sample to women who head a household or are married to its
head makes the young gradient steeper, not flatter.

\begin{table}[t]
\centering
\caption{Children ever born by family income (CPS, 2024)}
\label{tab:emp-fertility-income}
\begin{tabular}{lcccc}
\toprule
 & Bottom third & Middle third & Top third & Top minus bottom \\
\midrule
Women aged 24--26 & 0.596 & 0.448 & 0.226 & $-0.369$ \ $(0.052)$ \\
Women aged 40--44 & 1.784 & 1.702 & 1.733 & $-0.051$ \ $(0.050)$ \\
\bottomrule
\end{tabular}
\par\smallskip
\begin{minipage}{0.92\linewidth}
\footnotesize\textit{Notes:} June 2024 CPS Fertility Supplement. Mean
children ever born, capped at three, weighted by the person weight. Thirds
are equal-weight groups of family money income in the previous twelve months.
Approximate standard errors in parentheses, clustered by household. 1,673
women aged 24--26 and 3,278 aged 40--44.
% [DATA] output/model/credit_mechanism_20261004/measurement/tables.json.
\end{minipage}
\end{table}

\subsection{Large homes are owned}

Figure~\ref{fig:tenure-by-size} divides the homes of each size in the 2023
American Housing Survey (AHS) by tenure. Of the homes with four or more
bedrooms, 87 percent are owner-occupied and 9 percent are rented; of the homes
with at most one bedroom, 13 percent are owner-occupied and 74 percent are
rented. Measured in rooms, the unit used in the rest of the paper, 43.0
percent of owner-occupied homes have seven or more rooms, against 7.5 percent
of rented homes. The gap persists among families: in households with children,
owners occupy 6.8 rooms on average and renters 5.2. Nearly half of renter households
with children, 48.7 percent, live in homes with at most two bedrooms,
against 7.4 percent of owner households with children.
% [DATA] AHS 2023 ahs_children_space_mismatch.csv (code/data/ahs_supply_snapshot/
% output_ahs_family_unit_menu_national/). These shares describe the
homes households occupy, not the homes offered for rent. They show that, in
the existing stock, a family that wants a large dwelling typically buys one.
% [DATA] AHS 2023 national PUF; code/data/ahs_supply_snapshot/ (sandbox
% FIGURES.md maps the figure to its source).

\begin{figure}[t]
\centering
\includegraphics[width=0.62\linewidth]{ahs_tenure_share_within_size}
\caption{Tenure of the housing stock by number of bedrooms (AHS 2023)}
\label{fig:tenure-by-size}
\par\smallskip
\begin{minipage}{0.92\linewidth}
\footnotesize\textit{Notes:} Each bar divides the housing units with the
stated number of bedrooms into owner-occupied units, units rented for cash,
and vacant units. Units occupied without payment of cash rent and units whose
occupants have a usual residence elsewhere are excluded. Weighted by the
housing-unit weight. Source: American Housing Survey 2023, national
public-use file.
\end{minipage}
\end{figure}

\subsection{Who holds the large homes}

Most large owner-occupied homes are held by households without children at
home, and the largest group among them is old. Figure~\ref{fig:large-homes}
divides the owner-occupied homes with six or more rooms in 42 large
metropolitan areas in the 2023 American Community Survey (ACS) among six groups
of households, defined by the age of the householder and by whether a child of
the householder younger than 18 lives in the home. Households aged 60 to 85
without children at home hold 40.0 percent of these homes, and households
aged 40 to 59 without children at home another 21.8 percent. Households aged
22 to 39 with children hold 10.6 percent. Across all ages, 68.3 percent of the
large owner-occupied homes house no child of the householder, close to the
72.9 percent of all households in these areas without a child at home. What
the figure shows is therefore where the large homes sit by age, not that
childless households hold more than their share. I call the concentration of
large owner-occupied homes among older households the intergenerational
allocation of housing. The figure does not show that older households would
sell at some price, or that young families would have more children if they
held these homes.
% [DATA] ACS 2023 one-year, 42 metros, householders 22-85, 289,109 households;
% numbers as in the Oct 4 sandbox paper (reviewed there).
\mocktodo{name the builder of the large-homes figure in the data appendix
(not found in code/ by the Oct 4 sandbox review); the sandbox calls this
"intergenerational misallocation", which the title uses; choose the term.}

\begin{figure}[t]
\centering
\includegraphics[width=0.9\linewidth]{large_homes_by_age_children_acs2023}
\caption{Who holds the large owner-occupied homes (ACS 2023)}
\label{fig:large-homes}
\par\smallskip
\begin{minipage}{0.92\linewidth}
\footnotesize\textit{Notes:} Each bar is the share of owner-occupied homes
with six or more rooms held by households in the group; the six shares sum to
100 percent. Groups are defined by the age of the householder and by whether
an own child of the householder younger than 18 lives in the household.
Householders aged 22 to 85 in 42 large metropolitan areas, weighted by the
household weight; 289,109 households.
\end{minipage}
\end{figure}

\subsection{Housing around the first birth}

The Panel Study of Income Dynamics (PSID) follows the same families from 1968
to 2019. I date each woman's first birth from her full history of biological
children and estimate event studies of her household's housing around it,
comparing first-time mothers with women who never have children, with the
estimator of Sun and Abraham (2021) and person and year
fixed effects. The reference window is two to three years before the birth,
so that moves made while expecting a child count as responses. The sample
keeps one woman per household and year who is its reference person or spouse.
% [DATA] mock 02_empirical_evidence.tex and data_appendix.tex; 63,338
% person-years, 3,655 first-time mothers, 1,654 childless women.

Households add about one room within four years of the first birth.
Figure~\ref{fig:emp-rooms-path} reports the path. The dwelling is 0.41 rooms
larger than in the reference window in the year of the birth, 0.82 rooms one
to two years later, and 1.03 rooms (standard error 0.07) three to four years
after it, on a base of 4.7 rooms. Fixing household position before the birth,
which removes a pre-trend that reflects who heads a household, the increase at
three to four years is 1.47 rooms (0.05). The added space comes with a move
into ownership, after which families move less (Figure~\ref{fig:emp-tenure-moves}).
The ownership rate is 11 percentage points higher in the birth window and 16
points higher three to four years later, on a base of 41 percent. The
probability of having moved since the previous interview falls by 16 to 21
points after the birth and stays lower for a decade, and among those who do
move, the share moving for more space doubles, from 12 to about 25 percent.
The ACS shows the same pattern at a point in time: among household heads aged
30 to 55, those whose oldest child is under four are 12.8 percentage points
more likely to own than heads without children.
% [DATA] ACS 2005-06 recent-parent gap 0.12761 (same as Section 4 target).

\begin{figure}[t]
\centering
\includegraphics[width=0.88\linewidth]{first_birth_rooms_psid}
\caption{Rooms around the first birth (PSID)}
\label{fig:emp-rooms-path}
\par\smallskip
\begin{minipage}{0.92\linewidth}
\footnotesize\textit{Notes:} Interaction-weighted event-study estimates in
two-year windows of event time, with 95 percent confidence intervals from
standard errors clustered by person. One woman per household-year who is the
reference person or spouse; 63,338 person-years, 3,655 first-time parents and
1,654 childless women. The hollow marker is the omitted reference window,
years $-3$ and $-2$. Mean rooms in the reference window: 4.70.
\end{minipage}
\end{figure}

\begin{figure}[t]
\centering
\includegraphics[width=0.92\linewidth]{first_birth_tenure_moves_psid}
\caption{Ownership and moving around the first birth (PSID)}
\label{fig:emp-tenure-moves}
\par\smallskip
\begin{minipage}{0.92\linewidth}
\footnotesize\textit{Notes:} Same sample, estimator and reference window as
Figure~\ref{fig:emp-rooms-path}. ``Moved for more space'' equals one if the
household moved since the previous interview and gave more space as the first
reason. Means in the reference window: ownership 0.41, moved 0.58.
\end{minipage}
\end{figure}

\subsection{Implications for the model}

These facts motivate the main features of the model. Because children arrive
early and need space, fertility and housing are chosen jointly over the life
cycle, with a minimum of space for parents. Because large homes are mostly
owned, the size of dwelling a family can occupy depends on its tenure; the
model limits the size of rental units, so a family that wants a large home
must buy it, with a down payment it accumulates by saving. Because the large
homes sit with older households, the price of space for young families depends
on the whole age distribution of households; the model is an
overlapping-generations economy in which the house price clears one market
shared by all ages.
